%% ============================================================
%%  Chapter 9 — Conclusions and Future Work
%%  Dissertation: Radiation-Responsive Materials and Electronic
%%  Systems for Advanced Radiation Detection and Monitoring
%% ============================================================

\chapter{Conclusions and Future Work}
\label{ch:conclusions}

%% ------------------------------------------------------------
\section{Summary of Research Contributions}
\label{sec:conc_summary}
%% ------------------------------------------------------------

This dissertation has investigated radiation-responsive materials,
detector devices, alternative sensing concepts, and radiation-tolerant
digital electronics for advanced radiation detection and monitoring.
The research was organized around four complementary research thrusts
aligned with a materials-to-systems hierarchy:
(1)~characterizing the radiation response of photovoltaic-architecture
halide perovskite solar cells as current-mode radiation monitors,
(2)~developing and demonstrating perovskite semiconductor radiation
detectors for direct-conversion operation across X-ray, gamma, and
fast-neutron modalities,
(3)~investigating thermal radiation sensing concepts as an alternative,
non-electronic monitoring pathway for nuclear environments, and
(4)~implementing and validating radiation-tolerant FPGA-based digital
instrumentation through two complementary benchmarks tested under fast
neutron irradiation.  The following
sections consolidate the principal findings, discuss their
implications for nuclear monitoring systems, and identify directions
for future research.


%% ------------------------------------------------------------
\section{Key Findings}
\label{sec:conc_findings}
%% ------------------------------------------------------------

\subsection{Perovskite Solar Cells as Current-Mode Radiation Monitors}

The work presented in Chapter~\ref{ch:CsFAPbI} demonstrated that
photovoltaic-architecture Cs$_{0.1}$FA$_{0.9}$PbI$_3$ perovskite
solar cells can function as sensitive current-mode monitors for
ionizing radiation spanning X-ray, gamma, and potentially neutron
fields.  Under $^{137}$Cs gamma irradiation at 44.16~rad\,hr$^{-1}$,
an induced current density of 2.19~nA\,cm$^{-2}$ was measured, and a
distinct step response was confirmed upon insertion and withdrawal of
the source.  Extended irradiation with $^{60}$Co at
355.61~rad\,hr$^{-1}$ yielded 3.56~nA\,cm$^{-2}$, confirming that
the current-mode detection approach scales across a dose-rate range
spanning nearly an order of magnitude.  A linear degradation rate of
approximately 0.74~nA\,rad$^{-1}$ was observed during six-day
continuous $^{137}$Cs irradiation, attributable to cumulative defect
creation and glass darkening in the substrate.

The device bandwidth was found to be limited by the measurement
electronics rather than the intrinsic perovskite photoresponse, with
rise times of 33--84~ns depending on amplifier configuration.  Under
X-ray irradiation, the device produced measurable short-circuit
current proportional to beam intensity, although a characteristic slow
current decline was observed consistent with ion-migration field
screening.  Neutron response at the OSURR Fast Beam Facility was
statistically inconclusive for a single module, and the
signal-to-noise analysis indicated that scaling the detector array area
by a factor of ten to twenty would be necessary to resolve the
neutron-induced current above the combined gamma and electronic noise
backgrounds.


\subsection{Perovskite Semiconductor Radiation Detectors}

Chapter~\ref{ch:perovskite_detectors} extended the investigation to
direct-conversion semiconductor radiation detectors spanning three
radiation types.

For X-ray detection, the MAPbI$_3$-hBN polycrystalline device
exhibited a monotonically increasing sensitivity with applied bias,
reaching 0.411~$\mu$C\,(Gy\,cm$^2$)$^{-1}$ at 20~V, consistent with
Hecht equation predictions for improving charge collection at higher
fields.  Vacuum testing confirmed that the removal of ambient moisture
enhanced both sensitivity and stability, directly linking performance
degradation to moisture-accelerated ion migration.  The 2D perovskite
(PEA$_2$MA$_2$Pb$_3$I$_{10}$) detector with 5F-PEAI surface
treatment demonstrated that this moisture-accelerated instability can
be effectively mitigated: the fluorinated barrier layer raised the
breakdown voltage from $-10$~V to $-16$~V at 40\% relative humidity
and enabled X-ray sensitivities approaching
$\sim$1000~$\mu$C\,(Gy$_{\mathrm{air}}$\,cm$^2$)$^{-1}$ at 4~V bias,
among the highest reported for polycrystalline perovskite detectors of
comparable thickness.

For gamma-ray detection, three bromide perovskite single-crystal
systems were evaluated.  The CsPbBr$_3$ device produced a
statistically significant count-rate response to a $^{137}$Cs source
(17.45~cps vs.\ 9.11~cps background) at 1000~V bias, but high and
unstable leakage current (315--710~nA) prevented resolution of the
662-keV photopeak.  The FAPbBr$_3$ device achieved substantially lower
operating leakage ($5 \pm 2.5$~nA maximum), sustained stable
high-voltage operation for over five hours, and produced spectral
features attributable to both $^{137}$Cs and $^{57}$Co in a joint
acquisition at 1600~V bias, representing a proof-of-concept
demonstration of gamma spectroscopy in a solution-grown halide
perovskite crystal.  The MAPbBr$_3$ device failed to produce
functional Schottky contacts under the deposition conditions used,
motivating a revised contact deposition protocol at reduced sputter
power.

For fast neutron detection, MHyPbCl$_3$ crystals were characterised as
the first organic--metal halide perovskite semiconductor evaluated in a
reactor fast neutron beam.  The material combines a high hydrogen
density ($\sim 4.0 \times 10^{22}$~cm$^{-3}$, comparable to stilbene)
with an exceptional resistivity of
$4.43 \times 10^{11}$~$\Omega$\,cm and a mobility-lifetime product of
$9.1 \times 10^{-3}$~cm$^2$\,V$^{-1}$.  Waveform data acquired at
the OSURR Fast Beam Facility confirmed radiation-induced pulse
activity distinct from the beam background, and pulse-height analysis
produced spectra consistent with the expected broad proton-recoil
energy deposition distribution from a poly-energetic fast neutron beam.
Neutron--gamma discrimination remains a challenge due to the large
Pb content, and data acquisition limitations prevented collection of
sufficient statistics for quantitative efficiency analysis.

Across all three detection modes, ion migration under applied electric
fields emerged as the central performance-limiting mechanism.  The
correlation of device instability with ambient humidity, confirmed by
vacuum tests on MAPbI$_3$-hBN and by the systematic
humidity-dependence study on 2D perovskite detectors, identifies
surface passivation, such as the 5F-PEAI treatment, as the most
promising near-term strategy for extending the operational stability
of perovskite radiation detectors.


\subsection{Thermal Radiation Sensing Concepts}

Chapter~\ref{ch:thermal_sensing} presented the simulation and analysis
work performed in support of the Optical Fiber-Based Gamma Thermometer
(OFBGT) and the Optical Fiber-Based Neutron Dosimeter (OFBND).  Two
modeling approaches, a simplified Python model and an iterative MATLAB
model with temperature-dependent thermal conductivities, yielded
consistent predictions, forecasting a maximum temperature differential
of approximately 27--30$\,^{\circ}$C for the OFBGT in the CIF and
155--159$\,^{\circ}$C for the OFBND in the AIF of the OSURR.

Experimental testing of the OFBGT revealed a measured temperature
differential of approximately 17$\,^{\circ}$C, falling short of the
approximately 25$\,^{\circ}$C threshold required for reliable
operation.  The discrepancy was attributed to the omission of radiative
heat transfer in the original simulations; when radiative heat transfer
was included retrospectively, the model predictions agreed closely with
the experimental observations.  The Modulation Transfer Function
analysis confirmed that both sensors possess adequate axial resolution
for the spatial frequencies expected in the OSURR flux profiles.  The
OFBND has not yet been fabricated or tested, but the simulation results
indicate a favorable signal strength contingent on the inclusion of
radiative heat transfer in future design iterations.


\subsection{Digital Electronics Verification for Radiation
Environments}

Chapters~\ref{ch:digital_electronics} and~\ref{chap:fpga_benchmark}
addressed the design and experimental validation of two complementary
FPGA-based test structures under fast neutron irradiation at the OSURR
Fast Neutron Beam Facility.

The Flip-Flop Shift Register (FFSR) DMR test structure
(Chapter~\ref{ch:digital_electronics}) was deployed across five
Spartan-7 FPGAs during a two-day irradiation campaign.  During
Day~1 ($\sim$6--7$\times 10^9$~n/cm$^2$ fast fluence, gamma
shutter closed), zero DMR-detected SEUs were observed across all
five devices, although one device (FFSR1) exhibited a UART output
format change consistent with a configuration memory SEU that the
DMR counter was unable to detect.  During Day~2
($\sim$8--9$\times 10^{11}$~n/cm$^2$, gamma shutter open), two
devices experienced early functional dropouts consistent with
radiation-induced latch-up, one device experienced a severe multi-SEU
cluster event producing 2,048 detected upsets together with a second
configuration memory SEU affecting the UART baud rate, and PSU
current monitoring showed a sustained $\sim$12.5\% increase
consistent with TID-induced leakage current accumulation.

The fixed-point thermal control benchmark
(Chapter~\ref{chap:fpga_benchmark}) provided a complementary
assessment by exercising a closed-loop I-PD controller, forward-Euler
thermal plant emulation, nonlinear valve mapping, and DMR comparison
within a single synthesizable Verilog design.  Comparative analysis
of two telemetry acquisitions revealed a smooth exponential decay of
all plant state variables to zero over approximately 2,500~s, with
zero DMR mismatch detection throughout.  Bit-exact simulation
confirmed that the decay could not arise from firmware numerical
drift.  The failure signature was consistent with a configuration
memory (CRAM) SEU affecting arithmetic or routing logic shared by
both DMR chains, demonstrating a fundamental limitation of dual
modular redundancy when both redundant chains share common FPGA
configuration memory.

Together, these two independent observations, one from a
shift-register test structure and one from a closed-loop control
system, establish that configuration memory SEUs are a primary
concern for SRAM-based FPGAs at the fluence levels encountered in
nuclear instrumentation environments, and that effective mitigation
requires configuration memory scrubbing, triple modular redundancy
with physical placement constraints, or periodic functional
self-testing against known reference trajectories.


%% ------------------------------------------------------------
\section{Implications for Nuclear Monitoring Systems}
\label{sec:conc_implications}
%% ------------------------------------------------------------

The findings of this dissertation carry several implications for the
design and deployment of integrated radiation monitoring systems in
nuclear environments.

First, halide perovskite materials have been shown to produce
measurable electrical responses to X-rays, gamma rays, and fast
neutrons across a range of device architectures.  The
current-mode photovoltaic approach and the direct-conversion
semiconductor approach represent complementary detection
philosophies, the former offering large-area, zero-bias sensing
suitable for integrated dose monitoring and prompt detonation
forensics, and the latter offering the possibility of pulse-by-pulse
energy spectroscopy at the cost of applied high voltage and the
attendant ion migration challenges.  The demonstration of gamma
spectroscopy in a solution-grown FAPbBr$_3$ crystal and of fast
neutron sensitivity in MHyPbCl$_3$ establishes that perovskite
materials can address multiple radiation types within a common
materials platform, a capability not available from conventional
compound semiconductors without resorting to separate detector
materials for each radiation species.

Second, the thermal sensing concepts explored for the OFBGT and OFBND
represent an alternative, non-electronic sensing pathway that avoids
the radiation susceptibility of semiconductor detectors entirely.
Although the experimental results highlight the importance of
accounting for radiative heat transfer in the design, the
fiber-optic sensing modality is inherently radiation-hard and offers
distributed measurement capability along the entire sensor length,
making it well suited for spatially resolved monitoring of flux
profiles inside reactor cores.

Third, the FPGA verification results provide quantitative evidence
that SRAM-based FPGAs are susceptible to configuration memory upsets
at moderate fast neutron fluences and to latch-up and TID degradation
at higher fluence levels.  The inability of DMR alone to detect
configuration memory upsets that affect both redundant chains
identically underscores the necessity of a defense-in-depth strategy
combining TMR with physical separation, configuration scrubbing via
ICAP/PCAP readback, and periodic functional self-testing for any
FPGA-based instrumentation deployed in a sustained radiation
environment.

Fourth, the hierarchy investigated in this dissertation, from
radiation-responsive materials through detector devices, alternative
sensing approaches, and digital electronics
instrumentation, reflects the functional layers that must be
addressed in any complete radiation monitoring system.  The results
demonstrate that no single technology addresses all requirements:
perovskite detectors offer material sensitivity but face ion migration
stability challenges; thermal sensors offer radiation hardness but
sacrifice time resolution and spectroscopic capability; and FPGA-based
electronics offer reconfigurability but require dedicated radiation
mitigation strategies.  An integrated monitoring architecture would
therefore draw on each technology for the function to which it is best
suited, with perovskite or conventional semiconductor detectors
providing the primary sensing element, thermal sensors providing
complementary dose verification, and radiation-tolerant FPGAs
providing the digital signal processing and data acquisition backbone.


%% ------------------------------------------------------------
\section{Future Research Directions}
\label{sec:conc_future}
%% ------------------------------------------------------------

The work presented in this dissertation opens several avenues for
continued investigation.

\subsection{Perovskite Materials and Devices}

The central limitation identified across all perovskite detector
systems was ion migration under applied electric fields, exacerbated
by ambient humidity.  Future work should pursue the optimization of
surface passivation strategies, building on the 5F-PEAI treatment
demonstrated in Chapter~\ref{ch:perovskite_detectors}, with the goal
of enabling sustained high-bias operation under ambient conditions.
Encapsulation methodologies that maintain the low-humidity performance
observed in vacuum testing, without requiring a permanent vacuum
environment, would represent a significant step toward practical
deployment.

For gamma spectroscopy, the FAPbBr$_3$ proof-of-concept motivates the
growth of larger single crystals with improved surface preparation and
contact deposition.  The revised low-power sputtering protocol
identified as necessary for MAPbBr$_3$ should be evaluated
systematically across all bromide compositions.  Achieving
spectroscopic energy resolution comparable to that of established CZT
detectors remains a long-term goal that will require simultaneous
advances in crystal quality, contact metallurgy, and electronic noise
reduction.

For fast neutron detection, the MHyPbCl$_3$ results warrant follow-up
experiments with a dedicated digital pulse-processing system to replace
the oscilloscope-based waveform acquisition that limited the
statistical power of the present measurements.  Neutron--gamma
discrimination techniques, including pulse-shape discrimination
algorithms, should be investigated to assess whether the organic
cation's hydrogen content can be exploited for fast neutron detection
while rejecting the gamma background inherent to the lead-containing
lattice.

For the current-mode photovoltaic approach, the signal-to-noise
analysis of Chapter~\ref{ch:CsFAPbI} identified array area scaling as
the path to resolving neutron-induced currents.  Fabrication and
testing of multi-module perovskite solar cell arrays, with parallel
development of low-noise current amplification electronics, would
directly test this prediction and move the concept toward the
large-area, real-time monitoring applications envisioned for
post-detonation forensics.

\subsection{Thermal Radiation Sensing}

The OFBGT results highlight the critical importance of including
radiative heat transfer in design-stage simulations.  Future
iterations of the sensor should incorporate radiative corrections
from the outset and explore design modifications, such as reduced
gas-gap dimensions or alternative fill gases, to increase the
temperature differential above the operational threshold.  The OFBND
concept remains to be validated experimentally, and fabrication and
in-reactor testing of the borosilicate glass thermal mass should be
prioritized as a near-term objective.

\subsection{FPGA-Based Instrumentation}

The DMR blind spot identified in this work, configuration memory
upsets that corrupt both redundant chains identically, motivates the
development and irradiation testing of a TMR implementation with
explicit physical placement constraints ensuring that each voter group
occupies distinct configuration frames.  Integration of ICAP-based
configuration scrubbing with the FFSR or thermal control benchmark
would allow quantitative evaluation of scrubbing effectiveness against
the CRAM upset rates measured in this campaign.  Additionally, the
implementation and testing of automated latch-up detection and
power-cycling recovery circuitry is needed to address the functional
dropouts observed during the high-fluence Day~2 irradiation.  The
fixed-point thermal control benchmark, with its deterministic state
evolution and 100~Hz telemetry, is particularly well suited as a test
vehicle for evaluating these mitigation strategies in future
irradiation campaigns.

\subsection{System Integration}

Looking further ahead, the integration of perovskite radiation sensors
with radiation-tolerant FPGA-based readout electronics into a single
prototype monitoring system would represent a natural convergence of
the materials, device, and electronics work streams pursued in this
dissertation.  Such a prototype would provide a platform for
evaluating system-level performance metrics, including detection
sensitivity, energy resolution, long-term stability, and radiation
tolerance of the complete signal chain, under realistic nuclear
environment conditions.